\bookchapter{第二十五章：二手手機處理與數字遺產}

\chapterintro{嘿，諸位，經前二十餘章之“修煉”，信君已自一“數字小白”成長爲一名相當從容之手機使用者矣。然科技生活總有輪迴——新歡終會變舊愛。當君準備令老夥計“退役”時，君知何以體面、安全地告別乎？}

\booksection{一、舊手機之“安樂死”：安全出售或轉贈指南}

核心原則：數據清理非“刪除”，而乃“覆蓋”。許多人以爲以照片拖進回收站、清空相冊遂萬事大吉，結果手機轉手後，新主人用數據恢復軟件一掃，老照片、聊天記錄、身份證照片全回來矣——此種“社死現場”每年皆於上演。

\begin{enumerate}
  \item \textbf{自我"大掃除"}：完成終一次完整備份，退出所有賬號（iCloud、華爲/小米/三星賬號、Google賬號），關閉"查尋吾之手機"功能——否則下一位使用者將爲激活鎖攔於門外。
  \item \textbf{執行終極"格式化"}：現代手機之"恢復出廠設置"通常包含加密擦除——它銷燬之乃加密密鑰本身，密鑰一沒，整盤數據瞬間變爲亂碼。
  \item \textbf{最後之"障眼法"（僅老設備需要）}：現代手機（iOS 8+、Android 6+）已默認全盤加密（FBE），恢復出廠設置將銷燬加密密鑰，數據理論上不可恢復，\textbf{無需此步驟}。僅對極老舊之非加密設備（2015年前），方須於恢復出廠設置後用垃圾文件填滿存儲，再執行一次恢復出廠設置。
\end{enumerate}

\textbf{各品牌數據擦除路徑速查表}：儘管原理相通，但每個品牌以“恢復出廠設置”藏之位置略有不同，關鍵差異如下。

\begin{center}
\begin{tabularx}{\linewidth}{|X|X|X|}
\hline
\textbf{品牌} & \textbf{擦除入口} & \textbf{關鍵差異與留意事項} \tabularnewline
\hline
華爲 & 設置 > 系統與更新 > 重置 > 恢復出廠設置 & 需先退出華爲賬號並關閉"查尋吾之手機"，否則將觸發激活鎖；純淨模式下或須密碼驗證。 \tabularnewline
小米 / Redmi & 設置 > 吾之設備 > 恢復出廠設置 & 底部有"格式化虛擬SD卡"與"清除全部內容與設置"選項，二者皆要勾選；退出小米賬號前先關閉"查尋設備"。 \tabularnewline
OPPO / 一加 & 設置 > 其他設置 > 還原手機 > 徹底清除全部數據 & ColorOS 提供"徹底清除"模式，將多次覆蓋存儲；需輸入鎖屏密碼確認。 \tabularnewline
vivo & 設置 > 系統管理 > 備份與重置 > 清除所有數據 & 勾選"格式化手機存儲"方將清理SD卡區；需退出vivo賬號，關閉"查尋手機"。 \tabularnewline
三星 & 設置 > 通用管理 > 重置 > 恢復出廠設置 & 需先移除Samsung賬號並關閉"查尋吾之手機"；部分機型有"自動覆蓋"選項，建議開啓。 \tabularnewline
iPhone & 設置 > 通用 > 傳輸或還原iPhone > 抹掉所有內容與設置 & 須先退出Apple ID，否則下一位用戶爲iCloud激活鎖攔住；iOS 15+支持"退出時保留iCloud數據"以備後查。 \tabularnewline
\hline
\end{tabularx}
\end{center}

\begin{tcbwarn}{最容易被忽略之三件事}
\begin{itemize}
  \item \textbf{未退出賬號}：手機賣賜下家後爲遠程鎖定，自己亦下不來，終只能持身份證去官方售後解鎖。
  \item \textbf{SD卡未取出}：恢復出廠設置未能格式化外置SD卡里之照片，記得手動取出或單獨格式化。
  \item \textbf{U盤/OTG外設未斷開}：擦除過程中若連着外設，或誤清外設數據，亦或導致擦除中斷。
\end{itemize}

\end{tcbwarn}

\booksection{二、數字遺產：君之社交靈魂何以安放？}

手機裏之數據，早遂非“虛擬”之矣。它承載着汝之資金、社交關係、創作成果，乃至汝之“數字身份”。\textbf{一個人去世後，此資產該何以處理？}此不再乃一個抽象問題，而乃越來越多家庭真實面臨之難題——親屬不知道密碼、平臺不使繼承、加密貨幣永遠丟失。

\begin{tcbinfo}{法律依據：個人信息保護法第49條}
2021年施行之《中華人民共和國個人信息保護法》第49條明確規定：\textbf{自然人死亡之，其近親屬爲了自身之合法、正當利益，可對死者相關個人信息行使本章規定之查閱、複製、更正、刪除等權利；死者生前另有安排之除外。}

此意味着：

\begin{itemize}
  \item \textbf{近親屬有法律依據}向平臺申請處理死者之賬號與個人信息。
  \item "合法、正當利益"通常包括辦理遺產繼承、關閉賬號、刪除敏感隱私等。
  \item 死者生前若有遺囑或協議另行安排（如指定遺產聯繫人），自其約定。
  \item 平臺不可以"用戶隱私"爲由一概拒絕配合，但具體流程各家不同。
\end{itemize}

\end{tcbinfo}

\booksubsection{微信與支付寶：賬號不可繼承，但資金可以}

此兩個中國人最常用之App，規則各不相同：

\begin{itemize}
  \item \textbf{微信賬號}：依據《騰訊微信軟件許可及服務協議》，賬號所有權歸騰訊，用戶僅享有使用權，\textbf{賬號本身不能繼承、不能轉讓}。但賬號內之\textbf{零錢、零錢通資金}屬於個人財產，近親屬可憑死亡證明、親屬關係證明、繼承公證書向財付通申請提現或轉移。
  \item \textbf{支付寶}：賬號同樣不可直接繼承，但支付寶餘額、餘額寶、理財通等\textbf{資金屬於可繼承財產}。繼承人可經由支付寶"客服熱線"或線下提交死亡證明、繼承公證書辦理；螞蟻財富類理財產品需走基金公司流程。
  \item \textbf{聊天記錄}：法律上屬個人隱私，平臺未能主動提供予家屬。如涉及遺產線索，可通過訴訟或公證申請調取。
\end{itemize}

\begin{tcbwarn}{切勿爲之事}
\begin{itemize}
  \item \textbf{勿以死者賬號密碼直接交予家人"遂用"}——此違反平臺協議，資金爲凍結甚至觸發風控時舉證更難。
  \item \textbf{勿嘗試用死者手機收驗證碼"提現"}——或爲認定爲盜用，反而觸犯法律。
  \item \textbf{勿拖延}：部分賬號長期未登錄將爲自動註銷，資金尋回將更麻煩。
\end{itemize}

\end{tcbwarn}

\booksubsection{Apple“遺產聯繫人”：預先指定，事半功倍}

iOS 15.2 起蘋果推出“遺產聯繫人”功能，允許用戶提前指定信任之聯繫人，去世後該聯繫人可憑\textbf{死亡證明 + 遺產聯繫人訪問密鑰}申請訪問死者之iCloud數據（照片、備忘錄、通訊錄、備忘錄等）。操作步驟：

\begin{enumerate}
  \item 於iPhone上進入\textbf{[設置] > [汝之姓名] > [登錄與安全] > [遺產聯繫人]}。
  \item 點擊"添加遺產聯繫人"，使用面容ID確認身份。
  \item 選擇聯繫人（自通訊錄挑選），系統將生成一份\textbf{訪問密鑰}。
  \item 將訪問密鑰通過短信、郵件或打印紙質件交給對方，並提醒對方妥善保存。
  \item 可添加多個聯繫人作爲備份，避免單點失效。
\end{enumerate}

爲指定人無需Apple設備，憑密鑰與死亡證明即可於蘋果官網申請訪問。整個流程通常於7天內完成。

\booksubsection{Google“不活躍賬戶管理器”：自動處理身後事}

Google提供“不活躍賬戶管理器（Inactive Account Manager）”，能於用戶長期未登錄時自動執行預設操作——此乃目前最完善之數字遺產機制之一。

\begin{enumerate}
  \item 訪問\textbf{myaccount.google.com/inactive}，登錄Google賬號。
  \item 設置超時期限：3 / 6 / 12 / 18個月未登錄後觸發。
  \item 添加受信任聯繫人（最多10人），並設置每位可訪問何數據（Gmail、相冊、雲端硬盤等）。
  \item 可選：設置自動回覆郵件，告知來信者汝已無法處理。
  \item 最後選擇：到期後是\textbf{分享數據給聯繫人}，抑或\textbf{直接刪除賬戶}。
  \item 設置手機號作爲兜底驗證，避免誤觸發。
\end{enumerate}

\booksubsection{國內社交賬號：紀念化與註銷流程}

國內主流社交平臺對逝者賬號之處理分兩種思路：\textbf{“紀念化”（保留賬號使親友緬懷）}和\textbf{“註銷”（徹底刪除）}。家屬可憑死亡證明、親屬關係證明、申請人身份證向平臺客服申請。

\begin{itemize}
  \item \textbf{微博}：支持賬號"紀念化"，由親屬提供死亡證明申請，賬號將爲鎖定無法登錄但內容保留，主頁將顯示"該賬號已設置爲紀念賬號"。亦可申請註銷，需提供繼承公證。
  \item \textbf{抖音}：通過"吾 - 設置 - 反饋與幫助 - 賬號與安全 - 賬號註銷"或人工客服，提交死亡證明、親屬關係證明、申請人身份證，審覈通過後可註銷或轉爲紀念狀態。
  \item \textbf{B站}：客服渠道提交逝者證明及親屬證明，可申請賬號註銷或保留爲紀念頁；UP主之歷史稿件處理需結合版權聲明。
  \item \textbf{小紅書}：經由官方客服申訴，提交死亡證明、親屬關係證明，可申請註銷賬號並刪除內容。
\end{itemize}

\begin{tcbinfo}{關鍵提醒}
\begin{itemize}
  \item 各平臺規則將更新，辦理前先查閱官方客服最新指引。
  \item 所有證明文件建議先掃描備份再提交，原件妥善保管。
  \item 跨平臺一次性辦理時，可統一申請一份\textbf{繼承公證書}，省去反覆跑手續。
\end{itemize}

\end{tcbinfo}

\booksubsection{數字資產之繼承難題}

除矣社交賬號，越來越多之人持有\textbf{真正有價值之數字資產}，其繼承問題更加複雜：

\begin{itemize}
  \item \textbf{加密貨幣（比特幣、以太坊等）}：基於"私鑰即所有權"原則，\textbf{無私鑰遂等於資產消失}。繼承人無法通過法律文件向鏈上"申訴"。建議：將助記詞/私鑰分片存放於多個物理安全位置（如銀行保險箱+家裏保險櫃），並於遺囑中註明訪問方式。
  \item \textbf{交易所賬戶}：與銀行賬戶類似，可憑繼承公證書向交易所（如幣安、OKX）申請處置，但跨境交易所流程漫長。
  \item \textbf{域名}：屬於可繼承財產，繼承人於域名註冊商（如阿里雲、GoDaddy）提交繼承公證後可辦理過戶。
  \item \textbf{知識產權（原創文章、視頻、專利、軟件著作權）}：著作權中之財產權利可繼承（人身權不繼承），收益權可由繼承人享有至作者死後50年。
  \item \textbf{雲存儲與網盤}：內容數據歸屬個人，可繼承；但部分平臺協議規定賬號本身不可轉讓，需先下載再註銷。
\end{itemize}

\textbf{加密貨幣“消失”之真實代價}\par

據Chainalysis報告，全球已有\textbf{超過300萬枚比特幣}因私鑰丟失而永遠無法尋回，相當於現有總量之六分之一。若汝持有加密資產，請務必使至少一位汝絕對信任之人知道\textbf{恢復路徑}——否則此財富將永遠沉睡於區塊鏈上。

\booksubsection{一份“數字遺囑”模板}

傳統之財產遺囑管不了數字資產。建議汝每年更新一份“數字遺囑”附件，與正式遺囑一併存放。以下模板供參考：

\textbf{《個人數字資產處置附件》（模板）}\par

\textbf{第一部分：賬號清單}（僅錄賬號名與屬平臺，\textbf{不寫密碼}）

\begin{itemize}
  \item 社交類：微信（賬號xxx）、微博、抖音、B站、小紅書等
  \item 支付類：支付寶、微信支付、各銀行卡App
  \item 雲端類：iCloud、華爲雲、小米雲、Google賬號、微軟賬號
  \item 郵箱類：QQ郵箱、網易郵箱、Gmail等
  \item 其他：網盤、訂閱服務、遊戲賬號等
\end{itemize}

\textbf{第二部分：密碼託管方案}

\begin{itemize}
  \item 主密碼存放方式：使用密碼管理器（如Bitwarden、1Password），主密碼寫於紙質件中，封存於家中保險櫃。
  \item 恢復密鑰/助記詞：分片存放，至少兩處物理位置（如銀行保管箱+家中保險櫃）。
  \item 授權查詢人：姓名、身份證號、與本人關係、聯繫方式。
\end{itemize}

\textbf{第三部分：處置意願}

\begin{itemize}
  \item 各類賬號希望"紀念化"抑"註銷"。
  \item iCloud / Google 數據是否指定"遺產聯繫人"或"不活躍賬戶管理器"。
  \item 加密資產由誰繼承，私鑰獲取路徑說明。
  \item 版權類作品收益由誰享有。
\end{itemize}

\textbf{第四部分：聯繫人與平臺}

\begin{itemize}
  \item 信任執行人姓名、聯繫方式、備用聯繫人。
  \item 常駐平臺客服入口與所需證明材料清單。
  \item 本附件更新日期：\_\_\_\_年\_\_月\_\_日，下次複核日期：\_\_\_\_年\_\_月\_\_日。
\end{itemize}

\begin{tcbwarn}{關於"密碼"之法律提示}
將密碼直接寫進遺囑存在泄露風險，且遺囑公示後所有繼承人皆能看到。\textbf{更穩妥之作法乃：遺囑裏唯寫“賬號清單 + 處置意願 + 密碼託管位置”，真正之密碼通過密碼管理器主密碼單獨傳賜信任執行人}。此樣既符合法律要求，亦降低矣密碼泄露風險。

\end{tcbwarn}

\booksection{三、環保之最後一程}

一部手機裏含金、銀、銅、稀土等數十種金屬，隨意丟棄既污染土壤，亦浪費資源。正確之回收渠道有三類，各有側重：

\begin{center}
\begin{tabularx}{\linewidth}{|X|X|X|X|}
\hline
\textbf{渠道} & \textbf{回收價} & \textbf{數據安全} & \textbf{適用場景} \tabularnewline
\hline
官方以舊換新（華爲/蘋果/小米等） & 中等，常疊加補貼券 & 最高，官方有規範流程 & 方欲換同品牌新機時最劃算 \tabularnewline
三方回收平臺（愛回收、轉轉、回收寶等） & 略高，可比價 & 較高，有數據清除報告 & 欲直接變現、不綁定品牌 \tabularnewline
公益捐贈（如鄉村學校、愛心機構） & 無經濟回報 & 需自行徹底擦除 & 設備仍可用、欲作善事 \tabularnewline
發揮餘熱（自留改造） & 省下新設備錢 & 原機主本人續用，無風險 & 設備性能尚可、家有改造需求 \tabularnewline
\hline
\end{tabularx}
\end{center}

\begin{itemize}
  \item \textbf{官方回收/以舊換新}：最推薦之方式，流程規範，數據安全有保障。蘋果、華爲、小米等品牌官網皆有"以舊換新"入口，評估後可抵扣新機款。
  \item \textbf{正規回收平臺}：選擇信譽良好之大型回收平臺，下單後順豐上門取件，檢測後打款。務必選擇\textbf{提供"數據清除報告"之}平臺，並留憑證。
  \item \textbf{公益捐}：可聯繫當地鄉村學校、社區公益組織，舊手機可改裝成教學平板或老人機使用。\textbf{捐贈前務必執行完整之擦除流程，並卸除所有SIM/SD卡}。
  \item \textbf{發揮餘熱}：改成車載導航儀、智能家居中控、監控攝像頭、寶寶監視器、電子相框。一臺舊手機可再戰好幾年。
\end{itemize}

\begin{tcbinfo}{電子垃圾之真相}
一部手機之電路板裏含有\textbf{金、銀、鈀、銅}等貴金屬，每噸廢舊電路板含金量乃金礦石之數十倍。但若隨意丟棄，其中之鉛、汞、鎘將滲入土壤與地下水，污染持續數十年。\textbf{正確回收，既保護環境，也能使稀缺金屬循環再生。}

\end{tcbinfo}

\begin{tcbpractice}{實操：爲君之舊手機行一場"標準葬禮"}
\begin{enumerate}
  \item 終一次完整備份：照片上傳雲端、重要文件複製至電腦、聊天記錄導出存檔。
  \item 退出所有賬號：iCloud / 華爲 / 小米 / Google 賬號，關閉"查尋吾之手機"。
  \item 取出SIM卡與SD卡，斷開所有OTG外設。
  \item 執行加密恢復出廠設置，必要時手動勾選"格式化存儲"。
  \item （僅老設備）用垃圾文件填滿存儲，再恢復出廠設置一次。
  \item 擇它之下一站：官方以舊換新、正規回收平臺、公益捐贈，或轉贈家人。
  \item 同步更新汝之"數字遺產清單"，以已淘汰之設備自清單中劃掉。
\end{enumerate}

\end{tcbpractice}

\begin{tcbpractice}{本月任務清單}
\begin{enumerate}
  \item 爲汝目前之主力手機指定Apple"遺產聯繫人"或設置Google"不活躍賬戶管理器"。
  \item 用本節模板起草一份《個人數字資產處置附件》，與家人溝通存放位置。
  \item 檢查家中是否有2臺以上閒置舊手機，預約一次官方或正規平臺回收。
\end{enumerate}

\end{tcbpractice}
